\section{Introduction}

\subsection{The problem, stated plainly}

A set of six binary containers totalling 37\,GB arrived on a workstation from a
third party's machine. They were said to contain a 2025-era AAA video game. Three
questions had to be answered, in order, before a single frame could be rendered:

\begin{enumerate}
\item \textbf{Is the payload hostile?} The provenance was a friend's Windows PC,
      and the working hypothesis on arrival was that the files had come from an
      impersonation site rather than the genuine distributor. If any component
      were malicious, the correct action was deletion, not analysis.
\item \textbf{Is the payload protected?} Modern commercial anti-tamper systems
      insert a virtualising layer between the operating system and the program's
      own control flow. Such a layer, or the countermeasures against it, can be
      categorically incompatible with a compatibility layer such as
      Wine~\cite{wine}, in which case no amount of engineering effort would ever
      produce a running program.
\item \textbf{Can it be made to run?} The containers refused to open under their
      own installer, and the title targets Direct3D~12 with Shader Model 6.8,
      ray tracing, and enhanced barriers---the upper bound of what a translation
      layer can currently express.
\end{enumerate}

Each question turned out to have an answer that was cheaper and more decisive
than expected, and each was initially attacked with an unnecessarily expensive
method. That gap---between the method first reached for and the method that
actually settled the question---is the recurring subject of this paper.

\subsection{Why this is worth writing down}

Case studies of this shape are usually either forensic (``here is the malware
we found'') or compatibility-focused (``here is the Wine bug we fixed''). They
are rarely both, and the interesting results here live precisely at the seams:

\begin{itemize}
\item The malware question and the anti-tamper question were \emph{settled by the
      same single measurement}---an Authenticode signature validation---because a
      valid retail signature simultaneously proves nobody added code and proves
      nobody removed any (Section~\ref{sec:denuvo}).
\item The extraction failure that consumed two working sessions was not a
      compression, codec, or archive-integrity problem at all. It was a
      configuration file that a library looked for at one hard-coded absolute
      path, and the ``unsupported compression method'' error it produced was
      technically accurate and completely misleading (Section~\ref{sec:extract}).
\item The final graphics defect presented as neither a crash nor a hang.
      Nine threads were blocked inside modal dialogs that the window system
      created and never painted. The program was, in a strict sense, waiting for
      the user to click a button that had never been drawn on any display
      (Section~\ref{sec:deadlock}).
\item The fix, once found, was to replace a single 5\,MB shared library. It was
      found only after everything sophisticated had failed, because the component
      in question had been an unexamined constant throughout
      (Section~\ref{sec:fix}).
\end{itemize}

\subsection{Contributions}

\begin{enumerate}
\item \textbf{A signature-first argument for anti-tamper triage.} We show that
      Authenticode validation over current bytes is a sufficient and inexpensive
      proof of anti-tamper absence, and that entropy profiling, section-table
      analysis, and import-table analysis serve as independent corroboration
      rather than as the primary instrument (Section~\ref{sec:denuvo}).
\item \textbf{A characterisation of vendor-divergent shader translation.} We
      report, to our knowledge for the first time in the literature, a
      quantitative diff of the SPIR-V emitted by a single D3D12-to-Vulkan
      translator for the same title on two GPU vendors: 4047 of 4430 modules
      share a content hash but differ in bytes, and one vendor-specific SPIR-V
      extension appears in 4039 modules on one card and none on the other
      (Section~\ref{sec:shaders}).
\item \textbf{A diagnostic technique for invisible modal deadlock.} Counting
      windows of Win32 class \verb|#32770| through the compatibility layer's own
      debugger converts an unfalsifiable ``it hangs'' into a scalar that a
      bisection harness can optimise (Sections~\ref{sec:deadlock},
      \ref{sec:harness}).
\item \textbf{A full negative-results table.} Twenty-three distinct hypotheses
      were tested and eliminated by measurement. We publish all of them, because
      the value of this work to a future reader is mostly in what they should
      \emph{not} spend an evening on (Sections~\ref{sec:shaders},
      \ref{sec:fix}).
\item \textbf{A silent-failure taxonomy} drawn from the eleven distinct instances
      encountered here in which a tool reported success, or reported an accurate
      but misdirecting error, while doing nothing useful
      (Section~\ref{sec:lessons}).
\end{enumerate}

\subsection{Scope, ethics, and what was not done}
\label{sec:ethics}

This work was performed on hardware owned by the author, on files already in the
author's possession, for two purposes: determining whether those files were
dangerous, and achieving interoperability with an operating system the vendor
does not support.

Three boundaries are worth stating explicitly, because a reader could otherwise
infer more than was done.

\textbf{No protection was circumvented by this work.} The game executable was
never modified; the central forensic finding is precisely that it is
\emph{byte-identical} to the signed retail build (Section~\ref{sec:denuvo}). The
Steamworks emulation component present in the archive was pre-existing, was
authored by a third party, and is \emph{characterised} here---its file sizes,
its loader hook, its configuration---but was neither written, improved, nor
analysed for the purpose of extending it. Nothing in this paper constitutes a
method for defeating a protection system.

\textbf{The contrasting title was abandoned, not attacked.} A second title on the
same machine does carry a commercial anti-tamper system, and its accompanying
countermeasure is a hypervisor that loads beneath the Windows kernel. We report
this only to explain why that title is categorically out of scope for a
compatibility layer: Wine implements no Windows kernel into which a driver can
load, and on Linux the hardware virtualisation extensions are already owned by
the host. That is an architectural impossibility argument, and it terminated the
investigation rather than directing it.

\textbf{The engineering contribution is entirely on the Linux side.} Every fix
described in Sections~\ref{sec:extract} through \ref{sec:fix}---the namespace
patch, the configuration path, the virtual desktop, the EGL vendor list, the
translation-layer substitution---modifies open-source components or their
configuration. None of them touches the game.

\subsection{Roadmap}

Section~\ref{sec:system} describes the hardware and software under test, which
matters more than usual here because a hybrid-graphics laptop is the source of
at least three of the defects. Sections~\ref{sec:malware} and \ref{sec:denuvo}
answer the provenance and protection questions. Section~\ref{sec:extract} covers
the two extraction failures and the standalone harness built to isolate them.
Section~\ref{sec:reassembly} describes reconstructing a 54\,GB installation
without its installer, and the a-priori verification technique that made every
step checkable in advance. Sections~\ref{sec:graphics1} through \ref{sec:fix}
are the graphics bring-up: reaching a swapchain, a three-session misdiagnosis of
a ``dead'' Vulkan driver, the invisible modal deadlock, the vendor shader
divergence, and the fix. Section~\ref{sec:lessons} generalises.
Section~\ref{sec:appendix} is an artifact and reproduction index.

Throughout, a claim is only stated as established if a specific measurement
supports it, and retracted claims are marked as such rather than deleted---three
of the working hypotheses in this investigation were confidently wrong, and the
manner of their failure is part of the result.
